% ===========================================================================
\section{Introducción}
\label{sec:intro}

\begin{consigna}[Trabajo final \textemdash{} Pipeline del procesador RISC-V]
\small\raggedright\setlength{\parskip}{4pt}
\textbf{Consigna.} Implementar el pipeline del procesador RISC-V (etapas IF, ID, EX, MEM y WB);
una Debug Unit que permita enviar y recibir información al procesador mediante protocolo UART, y
una interfaz para visualizar los datos e interactuar con la Debug Unit (CLI, GUI y/o TUI).

\textbf{Instrucciones a implementar.}
\begin{itemize}[nosep]
  \item R-type: \texttt{add, sub, sll, srl, sra, and, or, xor, slt, sltu}.
  \item I-type: \texttt{lb, lh, lw, lbu, lhu, addi, andi, ori, xori, slti, sltiu, slli, srli,
        srai, jalr}.
  \item S-type: \texttt{sb, sh, sw}. \quad B-type: \texttt{beq, bne}. \quad J-type:
        \texttt{jal}. \quad U-type: \texttt{lui}.
\end{itemize}

\textbf{Requerimientos.} El procesador debe poder programarse y reprogramarse con comandos por
UART, sin resintetizar. El clock no debe verse intervenido en ninguna parte del proyecto. Las
decisiones de diseño se documentan con su porqué.

\textbf{Debug Unit.} Se envían a la PC por la UART el contenido de los 32 registros, de los
latches intermedios y de la memoria de datos usada.

\textbf{Carga de programa.} El programa está escrito en ensamblador, con un mecanismo para
traducirlo a lenguaje máquina y una instrucción HALT. El sistema permite programar el procesador
escribiendo la memoria de programa, reprogramarlo de manera dinámica, y responde: ¿es necesario
vaciar la memoria?, ¿y los registros?, ¿se necesita vaciar el pipeline?, ¿y la memoria de
programa?

\textbf{Modos de operación.} Continuo (un comando ejecuta el programa hasta el final y se
muestran todos los valores) y paso a paso (cada comando ejecuta un ciclo de clock y se muestran
los valores). En ambos casos el pipeline debe quedar vacío al terminar. Responder: ¿qué sucede si
en la memoria no se encuentra una instrucción de parada?

\textbf{Clock.} En la integración: ¿cuál es el camino crítico del sistema?, ¿genera skew y qué
consecuencias tiene? De encontrarse skew, encontrar la frecuencia óptima, generar métricas con
las herramientas de Vivado y aplicar esa frecuencia.
\end{consigna}

El trabajo final de la materia es un procesador \textbf{RISC-V RV32I segmentado en cinco etapas}
(IF, ID, EX, MEM, WB), con unidad de cortocircuitos (\emph{forwarding}), detección de riesgos y
una \textbf{Debug Unit} que se comunica con la PC por UART. Está escrito en Verilog, se implementa
en una Basys~3 (Artix-7 \texttt{xc7a35tcpg236-1}) con Vivado~2024.2 y se maneja desde la PC con un
debugger de texto propio.

Desde la PC se puede cargar un programa en ensamblador sin volver a sintetizar, ejecutarlo de
corrido o de a un ciclo por comando, y ver en cada paso los 32 registros, los cuatro registros de
segmentación, la memoria de datos usada y qué instrucción hay en cada etapa.

\subsection{Resultados}

\begin{itemize}[nosep]
  \item \textbf{ISA}: el subset RV32I de la consigna (32 instrucciones) más \texttt{halt}
        (\texttt{ebreak}), sección~\ref{sec:isa}.
  \item \textbf{Pipeline}: cinco etapas con forwarding desde EX/MEM y MEM/WB, bypass en el banco
        de registros, un ciclo de \emph{stall} por carga-uso y saltos resueltos en EX con
        \emph{flush} de dos ciclos (secciones~\ref{sec:pipeline} y~\ref{sec:riesgos}).
  \item \textbf{Debug Unit}: protocolo de tramas con checksum, 14 comandos y dos máquinas de
        estado, sin tocar el reloj: el core avanza o se detiene con una señal
        \texttt{enable} (sección~\ref{sec:debug}).
  \item \textbf{Reloj}: \qty{92}{\mega\hertz} con el Clock Wizard, la máxima frecuencia de un
        barrido de 65 a \qty{100}{\mega\hertz} que cumple timing (slack $+0{,}026$~ns,
        sección~\ref{sec:timing}).
  \item \textbf{Verificación}: dos oráculos independientes (estado final escrito a mano y un
        modelo de referencia ciclo a ciclo), un banco que compara los cuatro registros de
        segmentación en cada ciclo, 39 mutantes del RTL y la prueba en la placa: 22 de 22
        programas con ejecución continua y paso a paso, y 32 de 32 de las variantes sin
        riesgos (sección~\ref{sec:verif}).
  \item \textbf{Herramientas}: ensamblador, modelo de referencia, debugger con interfaz de texto y
        scripts de regresión, timing y prueba en placa (secciones~\ref{sec:interfaz}
        y~\ref{sec:uso}).
\end{itemize}

\subsection{Repositorio}
\label{sec:repo}

\begin{nota}[Código, documentación y evidencia]
\centering\normalsize
\url{\Repo}\par\smallskip
\small Versión descripta en este informe: tag \href{\Repo/tree/v1.0}{\texttt{v1.0}}
\end{nota}

Este informe resume el trabajo; lo completo está en el repositorio: el RTL, los testbenches, los
programas de prueba, las herramientas de la PC, la especificación, las 45 notas de decisión y los
reportes de Vivado. Las rutas que aparecen en el texto con este formato, como
\repo{docs/spec/pipeline.md} o \repo{tools/debugger/}, son relativas a la raíz del repositorio.
Para empezar conviene leer:

\begin{itemize}[nosep]
  \item \href{\Repo\#readme}{\texttt{README.md}}: estado, estructura y cómo usar el procesador.
  \item \href{\Repo/tree/master/docs}{\texttt{docs/}}: especificación, decisiones, diagramas,
        timing y capturas de la placa.
  \item \href{\Repo/tree/master/sw/demo}{\texttt{sw/demo/}}: los programas de demostración y el
        guion para recorrerlos con el debugger.
  \item \href{\Repo/blob/master/docs/informe/presentacion.pdf}{\texttt{docs/informe/presentacion.pdf}}:
        la presentación del trabajo.
\end{itemize}

\subsection{Cómo está organizado el trabajo}

Este informe resume y enlaza la documentación del repositorio:

\begin{center}
\small
\begin{tabular}{@{}p{4.2cm} p{11.4cm}@{}}
  \toprule
  \textbf{Carpeta} & \textbf{Contenido} \\
  \midrule
  \repo{docs/spec/}        & Especificación escrita antes del código: ISA, pipeline, memorias, protocolo, interfaz y verificación \\
  \repo{docs/decisiones/}  & 45 notas de decisión de diseño, con las opciones y el porqué (anexo~\ref{anx:decisiones}) \\
  \repo{docs/timing/}      & Reportes de Vivado, camino crítico, skew y barrido de frecuencias \\
  \repo{rtl/}, \repo{fpga/}, \repo{ip/} & Diseño sintetizable, top de la placa y scripts de los IP cores \\
  \repo{tb/}               & Testbenches unitarios y el banco de pruebas del pipeline \\
  \repo{sw/}               & Programas de prueba en ensamblador con su resultado esperado \\
  \repo{tools/}            & Ensamblador, modelo de referencia y debugger \\
  \repo{scripts/}          & Proyecto de Vivado, regresión, timing y barrido \\
  \bottomrule
\end{tabular}
\end{center}

\paragraph{Metodología.} El pipeline se construyó \textbf{etapa por etapa}: cada etapa fue un
hito con sus módulos, sus testbenches unitarios y una integración verificada contra el modelo de
referencia antes de pasar a la siguiente. Los riesgos se agregaron recién con las cinco etapas
funcionando, porque involucran a varias a la vez. Cada tarea fue una issue de GitHub (se citan
como I-NN) con sus criterios de terminado, una rama y un pull request revisado por el otro
integrante. Antes de escribir un módulo se documentó su interfaz en \repo{docs/spec/}, y cada
decisión de diseño quedó en una nota numerada. La planificación por hitos está en
\repo{docs/planificacion.md}; el avance real, en las issues y los pull requests del repositorio.
